\begin{textsample}{NEG Dim 1 – 2019\_c5 – Score 16.00 – 2019\_c5/enade\_2019\_nutri\_acdc34d1d9.txt}  \label{ex:f1_neg_399}
CAPÍTULO 5 DISTRIBUIÇÃO DOS CONCEITOS 5.1 PANORAMA NACIONAL DA DISTRIBUIÇÃO DOS CONCEITOS O Conceito Enade é um indicador de qualidade que avalia os cursos por intermédio dos desempenhos dos estudantes a partir dos resultados do Enade.
Os valores possíveis vão de 1 ( pior situação ) a 5 ( melhor situação ).
Os cursos com apenas um ou sem nenhum concluinte participante não obtêm o Conceito Enade, ficando Sem Conceito ( SC ).
Na Tabela 5.1, são apresentadas a quantidade e a distribuição de cursos de Nutrição participantes do Enade/2019, por faixa de conceito e Grande Região.
Enfatiza -se, mais uma vez, que a diferença entre os cursos tabulados neste capítulo e no capítulo 2 corresponde aos cursos que não foram avaliados, ou seja, àqueles, em princípio, sem alunos concluintes que tivessem sido inscritos para a prova.
Observando -se os dados da Tabela 5.1, nota -se que, dos 423 cursos participantes, 156 ( 36,9 \% ) classificaram -se com conceito 3, o valor modal.
Esse foi também o conceito modal nas seguintes regiões : Sudeste ( 37,8 \% ), Sul ( 44,1 \% ) e Centro-Oeste ( 41,9 \% ).
O conceito 2 foi o segundo mais frequente em nível nacional ( 24,8 \%, correspondendo a 105 cursos ), e o conceito 4, o terceiro ( 24,6 \%, correspondendo a 104 cursos ).
Houve, ainda, 41 cursos ( 9,7 \% ) que receberam conceito 5, e oito cursos ( 1,9 \% ) que receberam conceito 1.
Dos 423 cursos de Nutrição, nove ( 2,1 \% ) ficaram Sem Conceito ( SC ).
A região Norte participou com 23 cursos ou 5,4 \% do total nacional.
Desses, 12 cursos receberam o conceito 2, o que equivale a 52,2 \% do total regional, além de ser o conceito modal na região.
Os conceitos 3 e 4 foram atribuídos a quatro cursos \textbf{cada} um.
Um curso recebeu conceito 1.
Dois cursos ficaram Sem Conceito ( SC ) e nenhum recebeu conceito 5.
A região Nordeste participou com 113 cursos ou 26,7 \% do total nacional.
Nessa região, 46 cursos ( 40,7 \% em termos regionais ) obtiveram conceito 2, o conceito modal para a região.
Os conceitos 3 e 4 foram atribuídos, respectivamente, a 38 e a 15 cursos ( 33,6 \% e 13,3 \%, respectivamente ).
Ainda nessa região, 12 cursos ( 10,6 \% ) foram avaliados com o conceito 5.
Um curso ficou com conceito 1 e um curso ficou Sem Conceito ( SC ).
Dos 188 cursos participantes da região Sudeste ( 44,4 \% do total nacional ), o conceito 3 foi o conceito modal, como já anteriormente assinalado, tendo sido obtido por 71 cursos ( 37,8 \% ).
O conceito 1 foi atribuído a cinco cursos ( 2,7 \% ), e o conceito 2, a 36 cursos ( 19,1 \% ).
Receberam os conceitos 4 e 5, respectivamente, 50 cursos ( 26,6 \% ) e 21 cursos ( 11,2 \% ).
Nessa região, cinco cursos ficaram Sem Conceito ( SC ).
Os 68 cursos da região Sul corresponderam a 16,1 \% do total nacional.
Como já indicado, houve predominância do conceito 3, atribuído a 44,1 \% dos cursos da região ( conceito modal ), ou seja, atribuído a 30 dos 67 cursos participantes na região Sul.
O conceito 2 foi atribuído a cinco cursos ( 7,4 \% ) e os conceitos 4 e 5, a 27 e cinco cursos, \textbf{cada} ( 39,7 \% e 7,4 \%, respectivamente ).
Um curso ficou Sem Conceito ( SC ) e nenhum recebeu conceito 1.
Já dos 31 cursos participantes na região Centro-Oeste ( 7,3 \% do total nacional ), 13 ( 41,9 \% em termos regionais ) receberam conceito 3, o conceito modal, como já fora destacado.
Seis cursos ( 19,4 \% ) obtiveram conceito 2, e outros oito ( 25,8 \% ), conceito 4.
Três cursos ( 9,7 \% ) receberam conceito 5, um curso ( 3,2 \% ) recebeu conceito 1.
Nenhum curso ficou Sem Conceito ( SC ).
No Gráfico 5.1, é apresentada a distribuição cumulativa do Conceito Enade segundo Grande Região.
Em um gráfico de distribuição cumulativa, em \textbf{cada} valor de abscissa ( eixo x ), o valor apresentado no eixo das ordenadas ( eixo y ) é o valor acumulado da variável ( conceito Enade ).
Observando a linha vertical pontilhada, que ajuda a visualizar o percentual de conceitos acumulados até o valor 3, ou seja, conceitos 1, 2 e 3, pode -se notar que, por exemplo, a região Norte ( linha azul ) apresenta 81 \% dos cursos com conceito 3 ou abaixo deste valor.
A região Sul ( linha rosa ), por outro lado, apresenta 52 \% dos cursos na mesma situação : com conceito 3 ou abaixo.
Destaca -se que em algumas regiões o conceito 5 pode não ser alcançado por nenhum curso, como ocorre na região Norte ( linha azul ), na qual todos os cursos têm conceito menor ou igual a 4, e a poligonal alcança 100 \% no conceito 4.
Em linhas gerais, em um gráfico de distribuição cumulativa, poligonais mais à esquerda representam Grandes Regiões com pior distribuição, e poligonais mais à direita, Grandes Regiões com melhores conceitos.
A região Norte ( linha azul ) apresenta os cursos com conceitos mais baixos, ao passo que as regiões Sudeste ( linha vermelha ) e Sul ( linha rosa ) apresentam as distribuições com valores maiores.
Cursos em IES situadas na região Centro-Oeste ( linha laranja ) apresentam uma situação intermediária. 5.2 CONCEITOS POR CATEGORIA ADMINISTRATIVA, POR MODALIDADE DE ENSINO E GRANDE REGIÃO Na Tabela 5.2, é apresentada a distribuição dos cursos de Nutrição participantes do Enade/2019, por Categoria Administrativa e Modalidade de ensino, de acordo com os conceitos por eles alcançados, segundo as Grandes Regiões brasileiras.
Dos 423 cursos participantes, 71 ( 16,8 \% ) eram ministrados em Instituições Públicas, e 352 ( 83,2 \% ), em Instituições Privadas.
Quanto à Modalidade de ensino, dois cursos eram a Distância, e 421, Presenciais.
De acordo com as informações da Tabela 5.2, em termos nacionais, entre Instituições Públicas e Privadas, observa -se que, dos 41 cursos avaliados com conceito 5, 35 eram oferecidos em IES Públicas, e seis, em IES Privadas.
Dos 71 cursos participantes de IES Públicas, o conceito 5 foi o valor modal, atribuído a 35 cursos ( 49,3 \% ).
Dos demais cursos de Instituições Públicas participantes, quatro cursos foram avaliados com conceito 2, 21 cursos ( 29,6 \% ) foram avaliados com conceito 4 e, como já comentado, 35 cursos ( 49,3 \% ), com conceito 5.
Outros 11 cursos receberam conceito 3.
Nenhum curso ficou Sem Conceito ( SC ) ou recebeu conceito 1.
Na rede privada, o conceito modal foi 3, atribuído a 145 cursos, o correspondente a 41,2 \% dos 352 cursos da categoria.
Dos demais cursos participantes, oito ( 2,3 \% ) receberam conceito 1 e 101 ( 28,7 \% ) conceito 2.
O conceito 4 foi atribuído a 83 cursos ( 23,6 \% ), e o conceito 5 a seis cursos ( 1,7 \% ).
Nesta Categoria Administrativa, nove cursos ( 2,6 \% ) ficaram Sem Conceito ( SC ).
Dos 421 cursos participantes da modalidade presencial, o conceito 3 foi o valor modal, atribuído a 154 cursos ( 36,6 \% ).
Dos demais cursos nesta modalidade de ensino, oito cursos ( 1,9 \% ) receberam conceito 1, 105 cursos ( 24,9 \% ) receberam conceito 2, 104 cursos ( 24,7 \% ), conceito 4, e 41 cursos ( 9,7 \% ), conceito 5.
Nesta modalidade de ensino, nove cursos ( 2,1 \% ) ficaram Sem Conceito ( SC ).
Aos dois únicos cursos a Distância informados atribuiu -se o conceito 3.
Na análise por região, observa -se que, na região Norte, as Instituições Privadas participaram com 19 cursos ( 83,2 \% do total regional ), dos quais dois ficaram Sem Conceito ( SC ) e nenhum recebeu o conceito 5.
O conceito modal para as Instituições Privadas na região foi 2, atribuído a 11 cursos, correspondendo a 57,9 \% dessa categoria na região.
Quatro cursos foram avaliados com conceito 3.
Os conceitos 1 e 4 foram atribuídos a um curso \textbf{cada}.
As Instituições Públicas participaram com quatro cursos na região Norte ( 17,4 \% do total regional ), aos quais foram atribuídos os conceitos 2 e 4.
Nenhum curso ficou Sem Conceito ( SC ) ou com os conceitos 1, 3 e 5 nessa combinação de categoria e região.
Todos os cursos oferecidos na região Norte eram presenciais.
Na região Nordeste, a rede privada concentrou 93 dos 113 cursos participantes, o equivalente a 82,3 \% do total da região.
Dentre os cursos oferecidos por IES Privadas no Nordeste, 46 cursos foram avaliados com conceito 2, o conceito modal.
Receberam conceito 3, 36 cursos.
Outros oito cursos receberam conceito 4 e os conceitos 1 e 5 foram atribuídos a um curso \textbf{cada}.
Nessa combinação de categoria e região, um curso ficou Sem Conceito ( SC ).
As Instituições Públicas da região Nordeste participaram com 20 cursos ( 17,7 \% do total da região ).
Desses, 11 obtiveram conceito 5, o valor modal, seguidos por outros sete cursos, que receberam conceito 4, e de dois cursos que receberam conceito 3.
Nenhum curso ficou Sem Conceito ( SC ) ou com os demais conceitos nessa região.
Todos os cursos oferecidos na região Nordeste eram presenciais.
Na região Sudeste, a proporção de cursos da rede privada, 83,5 \%, correspondendo a 157 dos 188 cursos participantes.
Nessa categoria e região, o conceito modal foi 3, atribuído a 67 cursos, e cinco cursos ficaram Sem Conceito ( SC ).
Os demais foram avaliados com conceito 1 ( cinco cursos ), conceito 2 ( 33 cursos ), conceito 4 ( 43 cursos ) e conceito 5 ( quatro cursos ).
Dos 31 cursos oferecidos por Instituições Públicas na região Sudeste, o conceito modal foi 5, atribuído a 17 cursos.
Sete cursos receberam conceito 4, quatro receberam conceito 3, e três cursos receberam conceito 2.
Nessa região, nenhum curso ficou Sem Conceito ( SC ) ou foi avaliado com o conceito 1.
Também na região Sudeste, a maioria dos cursos foi oferecida na modalidade Presencial : 187 cursos.
Para esses, a moda foi o conceito 3, com 70 cursos.
Os demais cursos presenciais da região foram enquadrados nos conceitos 1 ( cinco cursos ), 2 ( 36 cursos ), 4 ( 50 cursos ) e 5 ( 21 cursos ).
Cinco cursos ficaram Sem Conceito ( SC ).
Na região Sudeste, o curso oferecido na modalidade Educação a Distância obteve conceito 3.
As Instituições Privadas concentraram 57 dos 68 cursos participantes da região Sul, 83,8 \% do total regional.
Desses, 26 cursos obtiveram conceito 3, o conceito modal.
Nesta combinação de Categoria Administrativa e Grande Região, cinco cursos receberam conceito 2, outros 24, conceito 4, e um curso, conceito 5.
Nenhum curso recebeu conceito 1 e um curso ficou Sem Conceito ( SC ).
As Instituições Públicas na região Sul participaram com 11 cursos ( 16,2 \% dos cursos da região ), aos quais foram atribuídos os conceitos 3, 4 e 5.
Nenhum curso da região Sul oferecido por IES Pública ficou Sem Conceito ( SC ) ou recebeu os conceitos 1 e 2.
Todos os cursos oferecidos na região Sul eram presenciais.
Na região Centro-Oeste, 26 dos 31 cursos participantes eram de Instituições Privadas ( 83,9 \% em termos regionais ).
Desses, 12 receberam conceito 3, o conceito modal.
Dos demais cursos, sete receberam conceito 4, seis receberam conceito 2 e um recebeu conceito 1.
Nenhum curso ficou Sem Conceito ( SC ) ou recebeu conceito 5.
Quanto aos cinco cursos oferecidos por Instituições Públicas na região Centro-Oeste ( 16,1 \% do total regional ), três cursos apresentaram conceito 5, conceito modal.
Os conceitos 3 e 4 foram atribuídos a um curso, \textbf{cada}.
Nenhum curso ficou Sem Conceito ( SC ) ou recebeu os conceitos 1 e 2.
Todos os cursos oferecidos na região Centro-Oeste eram presenciais.
No Gráfico 5.2, apresenta -se a distribuição do Conceito Enade, segundo a Categoria Administrativa da IES.
Os cursos em IES Públicas ( linha azul ) apresentam uma poligonal mais à direita do que os em IES Privadas ( linha verde ) e, consequentemente, uma distribuição de Conceitos Enade com valores maiores. 5.3 CONCEITOS POR ORGANIZAÇÃO ACADÊMICA E GRANDE REGIÃO Na Tabela 5.3, encontra -se a distribuição dos conceitos atribuídos aos cursos participantes do Enade/2019, na Área de Nutrição, por Organização Acadêmica, segundo as Grandes Regiões brasileiras.
Dos 423 cursos de Nutrição participantes, 179 eram oferecidos em Universidades, 131 em Centros Universitários, 111 em Faculdades e dois em CEFET/IF.
Esta distribuição corresponde a, respectivamente, 42,3 \%, 31,0 \%, 26,2 \% e 0,5 \% dos cursos.
De acordo com os dados apresentados, de todos os 41 cursos avaliados com conceito 5, 39 eram vinculados a Universidades.
Os cursos vinculados a esse tipo de Organização Acadêmica tiveram o conceito 4 como conceito modal, atribuído a 61 cursos ( 34,1 \% ).
Os demais cursos vinculados a Universidades receberam os conceitos 1 ( quatro cursos ), 2 ( 21 cursos ), 3 ( 50 cursos ) e conceito 5 ( 39 cursos, como já mencionado ).
Ficaram Sem Conceito ( SC ), quatro cursos.
Dos cursos em Centros Universitários, o conceito modal foi 3, atribuído a 61 cursos ( 46,6 \% ).
Três cursos vinculados a esse tipo de Organização Acadêmica ficaram Sem Conceito ( SC ), e os outros cursos receberam os conceitos 1 ( dois cursos ), 2 ( 35 cursos ), 4 ( 29 cursos ) e 5 ( um curso ).
Dos 111 cursos mantidos por Faculdades, 49 ( 44,1 \% ) receberam conceito 2, o conceito modal.
O segundo conceito mais frequentemente atribuído aos cursos oferecidos por essa categoria de IES foi o conceito 3, recebido por 45 cursos.
Na sequência, 12 cursos obtiveram conceito 4, dois cursos, conceito 1, e um curso, conceito 5.
Dos cursos oferecidos em Faculdades, dois ficaram Sem Conceito ( SC ).
Dos dois cursos mantidos por CEFET/IF receberam conceito 4.
Considerando -se, separadamente, as regiões brasileiras, verifica -se que, na região Norte, as Universidades participaram com sete dos 23 cursos avaliados.
Três cursos receberam conceito 4, conceito modal, outros dois receberam conceito 2.
Os conceitos 1 e 3 foram atribuídos a um curso, \textbf{cada}.
Nenhum curso ficou Sem Conceito ( SC ) ou recebeu conceito 5.
Os Centros Universitários da região Norte foram representados por seis cursos, os quais obtiveram conceitos 2 ( três cursos ), 3 ( um curso ) e 4 ( um curso ).
Nenhum curso ficou com conceito 1 ou 5.
Um curso ficou Sem Conceito ( SC ).
As Faculdades participaram com dez cursos na região Norte, os quais receberam os conceitos 2 ( sete cursos, conceito modal ) e 3 ( dois cursos ) e um curso ficou Sem Conceito ( SC ).
Já os CEFET/IF não participaram com oferta de cursos.
Na região Nordeste, as Universidades participaram com 28 dos 113 cursos da Área de Nutrição oferecidos.
Nenhum desses ficou Sem Conceito ( SC ) ou recebeu conceito 1.
O conceito modal foi 5 ( 11 cursos ), seguido dos conceitos 3 e 4 ( sete cursos, \textbf{cada} ) e do conceito 2, atribuído a três cursos.
Os Centros Universitários contaram com 40 cursos participantes na região Nordeste, 19 dos quais receberam o conceito modal 2.
Para os demais cursos, foram atribuídos os conceitos 1 ( um curso ), 3 ( 17 cursos ) e 4 ( três cursos ).
Nenhum curso ficou Sem Conceito ( SC ) ou recebeu conceito 5.
As Faculdades foram representadas por 44 cursos na região Nordeste, 24 dos quais receberam conceito 2, o conceito modal.
O conceito 3 foi atribuído a 14 cursos, o conceito 4 a quatro cursos e o conceito 5 a um curso.
Nenhum curso recebeu conceito 1.
Um curso ficou Sem Conceito ( SC ).
Os CEFET/IF participaram com apenas um curso, que recebeu conceito 4.
Na região Sudeste, as Universidades concentraram 95 dos 188 cursos de Nutrição da região.
Dos cursos oferecidos em Universidades, nessa região, o conceito modal foi 4, atribuído a 30 cursos.
Os demais cursos receberam os conceitos 1 ( três cursos ), 2 ( 13 cursos ), 3 ( 26 cursos ) e 5 ( 20 cursos ).
Três cursos ficaram Sem Conceito ( SC ).
Os Centros Universitários participaram com 54 cursos na região Sudeste, dos quais 29 obtiveram o conceito modal 3.
O conceito 2 foi atribuído a dez cursos, o conceito 4, a 12 cursos e o conceito 5 a um curso.
Nenhum curso recebeu conceito 1.
Dois cursos ficaram Sem Conceito ( SC ), nessa combinação de organização acadêmica e região.
Já as Faculdades foram representadas por 38 cursos na região Sudeste, dos quais nenhum ficou Sem Conceito ( SC ) ou recebeu conceito 5.
O conceito modal 3 foi atribuído a 16 cursos.
Aos demais cursos foram atribuídos os conceitos 1 ( dois cursos ), 2 ( 13 cursos ) e 4 ( sete cursos ).
Os CEFET/IF participaram com apenas um curso, que recebeu conceito 4.
Dos 68 cursos da região Sul, 37 eram mantidos por Universidades, dos quais o conceito modal foi 4, atribuído a 16 cursos.
Os demais cursos receberam os conceitos 2 ( três cursos ), 3 ( 12 cursos ) e 5 ( cinco cursos ).
Nesse tipo de organização, um curso da região Sul ficou Sem Conceito ( SC ).
Nenhum curso recebeu conceito 1.
Dos 21 cursos participantes de Centros Universitários da região Sul, nenhum ficou Sem Conceito ( SC ) ou recebeu os conceitos 1 e 5.
Os conceitos 3 e 4 foram atribuídos a dez cursos, \textbf{cada}.
Um curso recebeu conceito 2.
Quanto aos dez cursos vinculados a Faculdades na região Sul, oito receberam o conceito modal 3.
Os conceitos 2 e 4 foram atribuídos a um curso, \textbf{cada}.
Nenhum curso ficou Sem Conceito ( SC ) ou recebeu os demais conceitos.
Já os CEFET/IF não participaram com oferta de cursos.
Na região Centro-Oeste, 12 dos 31 cursos oferecidos eram mantidos por Universidades.
Para este tipo de Organização Acadêmica, nenhum curso ficou Sem Conceito ( SC ) ou recebeu os conceitos 1 e 2.
O conceito modal foi 4 ( cinco cursos ), seguido do conceito 3 ( quatro cursos ) e 5 ( três cursos ).
Os Centros Universitários da região Centro-Oeste contaram com dez cursos.
Quatro cursos foram avaliados com conceito 3, o conceito modal.
Três cursos receberam conceito 4, dois receberam conceito 2 e um recebeu conceito 1.
Nenhum curso recebeu conceito 5 ou ficou Sem Conceito ( SC ).
Nesta região, as Faculdades ofereceram nove cursos, cinco dos quais ficaram com conceito 3 e quatro ficaram com conceito 2.
Nenhum curso ficou Sem Conceito ( SC ) ou recebeu os demais conceitos.
Já os CEFET/IF não participaram com oferta de cursos.
No Gráfico 5.3, apresenta -se a distribuição do Conceito Enade dos cursos de Nutrição, segundo a Organização Acadêmica da IES.
Os cursos em CEFET/IF ( linha laranja ) apresentam uma poligonal mais à direita do que as demais, denotando uma melhor distribuição de Conceitos Enade.
No outro extremo, com distribuição muito semelhante entre si, encontram -se os cursos mantidos por Centros Universitários ( linha verde ) e Faculdades ( linha vermelha ), apresentando as respectivas poligonais mais à esquerda e, consequentemente, distribuição de Conceitos Enade com menores valores.
As Universidades ( linha azul ) apresentam distribuições intermediárias.

% matched lemmas: cada
\end{textsample}
